\section*{अध्याय \ref{ch:g12:diversification-of-life} --- जीवधारियों का विविधीकरण}

\begin{solution}{exo:g12:diversification-of-life:1}
\omterm{def:g10:universal-dna:gene}{जीन} द्विगुणन और अपसरण; \omterm{prop:g11:antibiotic-resistance:variation}{क्षैतिज जीन अंतरण}; \omterm{prop:g12:diversification-of-life:polyploidy}{बहुगुणिता} के साथ संकरण; \omterm{prop:g12:diversification-of-life:development}{विकासीय
जीनों} के नियमन में बदलाव; \omterm{def:g12:diversification-of-life:symbiosis}{सहजीविता}; हस्तांतरित बरताव।
\end{solution}

\begin{solution}{exo:g12:diversification-of-life:2}
एक ही पूर्वज \omterm{def:g10:universal-dna:gene}{जीन} से लगातार द्विगुणनों द्वारा उतरे संबंधित \omterm{def:g10:universal-dna:gene}{जीनों} का एक
समुच्चय, जिसकी हर प्रति फिर \omterm{def:g11:mutations:mutation}{उत्परिवर्तन} से अपसरित होकर किसी अलग प्रकार्य पर
पहुँची (\omterm{def:g11:the-eye:photoreceptors}{ऑप्सिन}, ग्लोबिन)।
\end{solution}

\begin{solution}{exo:g12:diversification-of-life:3}
उसके \omterm{def:g10:universal-dna:information}{गुणसूत्रों} के पास, जिनका एक सेट हर जनक \omterm{def:g10:biodiversity-scales:species}{जाति} से आया है, \omterm{def:g12:meiosis-genetic-shuffling:meiosis}{अर्धसूत्री
विभाजन} पर जोड़ी बनाने को कोई समजात नहीं होता, इसलिए युग्मक असंतुलित होते
हैं। द्विगुणन हर \omterm{def:g11:cell-cycle-mitosis:chromosome}{गुणसूत्र} को एक समरूप साथी दे देता है: जोड़ी बनना और
\omterm{def:g12:meiosis-genetic-shuffling:meiosis}{अर्धसूत्री विभाजन} नियमित हो जाते हैं और वह पौधा उर्वर होता है।
\end{solution}

\begin{solution}{exo:g12:diversification-of-life:4}
उनका अपना वृत्ताकार \omterm{def:g10:universal-dna:information}{DNA}; जीवाणु प्रकार के \omterm{def:g10:cells-common-unit:organelle}{राइबोसोम}, जो जीवाणुरोधी
\omterm{def:g11:antibiotic-resistance:antibiotic}{प्रतिजैविकों} के प्रति संवेदनशील हैं; विभाजन से गुणन; एक दोहरी झिल्ली; और
जीवित जीवाणुओं के किसी समूह के सबसे पास पड़ते अनुक्रम।
\end{solution}

\begin{solution}{exo:g12:diversification-of-life:5}
वह बरताव जो समूह के दूसरे सदस्यों से सीखा जाता है और बिना \omterm{def:g10:universal-dna:gene}{जीनों} के आगे
बढ़ाया जाता है: कुछ चिंपैंज़ी समुदायों में पत्थरों से मेवे तोड़ना और
बाक़ियों में नहीं।
\end{solution}

\begin{solution}{exo:g12:diversification-of-life:6}
$14 + 7 = 21$ \omterm{def:g11:cell-cycle-mitosis:chromosome}{गुणसूत्र}: आइनकोर्न से 7 A, और एमर से 7 A तथा 7 B। सात जोड़े
(A के साथ A) बन सकते हैं; 7 B \omterm{def:g10:universal-dna:information}{गुणसूत्रों} के पास कोई साथी नहीं है और वे
यादृच्छिक ढंग से बँटते हैं, इसलिए युग्मक असंतुलित हैं: बंध्य।
\end{solution}

\begin{solution}{exo:g12:diversification-of-life:7}
$\alpha$--$\beta$ का बँटवारा (50\% समरूप) $\beta$--$\gamma$ के बँटवारे
(80\%) से पुराना है: पहले $\alpha$ शाखा बनाता है, फिर $\beta$ और $\gamma$
एक-दूसरे से अलग होते हैं।
\end{solution}

\begin{solution}{exo:g12:diversification-of-life:8}
पाद वाले \omterm{def:g10:universal-dna:gene}{जीन} मौजूद हैं पर वह संकेत, जो उन्हें भ्रूण के पार्श्व में चालू
करता है, दिया ही नहीं जाता, या दिया जाकर फिर रोक दिया जाता है; वह अंतर
अभिव्यक्ति के नियमन में है, उन \omterm{def:g10:universal-dna:gene}{जीनों} में नहीं।
\end{solution}

\begin{solution}{exo:g12:diversification-of-life:9}
भ्रूणीय चोंच में किसी वृद्धि संकेत की मात्रा और समय में: यानी नियामक अंतर।
नियमन के किसी बदलाव को किसी नियंत्रण अनुक्रम में केवल एक छोटा \omterm{def:g11:mutations:mutation}{उत्परिवर्तन}
चाहिए, जबकि किसी नए \omterm{def:g11:gene-expression:protein}{प्रोटीन} को बहुत-से चाहिए होते; इसलिए वह हज़ारों पीढ़ियों
के भीतर उठ सकता है और चुना जा सकता है।
\end{solution}

\begin{solution}{exo:g12:diversification-of-life:10}
उसका रूप, नंगी चट्टान बसाने की उसकी क्षमता, सूखे के प्रति उसका प्रतिरोध और
उसकी धीमी बढ़त किसी भी साझीदार के अकेले नहीं हैं; वे उस साझेदारी से उभरते
हैं। लाइकेन अपने ही गुणों वाला एक मिश्र जीव है।
\end{solution}

\begin{solution}{exo:g12:diversification-of-life:11}
हस्तांतरित बरताव: सामग्री दोनों ओर उपलब्ध है, वे आबादियाँ गहरे संबंध वाली
हैं, और वह अंतर उस नदी के पीछे चलता है --- यानी संपर्क में, और इसलिए सीखने
में, एक रुकावट, न कि \omterm{def:g10:universal-dna:gene}{जीनों} या संसाधनों में।
\end{solution}

\begin{solution}{exo:g12:diversification-of-life:12}
किसी विषाणु ने कुछ करोड़ साल पहले किसी नर-वानर पूर्वज की जनन \omterm{def:g10:cells-common-unit:cell}{कोशिका} के
जीनोम में अपना \omterm{def:g10:universal-dna:gene}{जीन} घुसा दिया; वह \omterm{def:g10:universal-dna:gene}{जीन} विरासत में मिलता गया और बचा रहा।
जाँच: वही \omterm{def:g10:universal-dna:gene}{जीन} सारे नर-वानरों में उसी गुणसूत्रीय जगह पर बैठा होना चाहिए और
बाक़ी स्तनधारियों में उस जगह से ग़ैरहाज़िर; और उसका अनुक्रम उस विषाणु के
अनुक्रम के सबसे पास होना चाहिए।
\end{solution}

\begin{solution}{exo:g12:diversification-of-life:13}
\omterm{def:g10:cells-common-unit:organelle}{सूत्रकणिका} \omterm{def:g10:cells-common-unit:organelle}{राइबोसोम} जीवाणु प्रकार के हैं और उन्हीं \omterm{def:g11:antibiotic-resistance:antibiotic}{प्रतिजैविकों} से बँधते
हैं: ऊँची ख़ुराकों पर वह दवा \omterm{def:g10:cells-common-unit:organelle}{सूत्रकणिकाओं} का \omterm{def:g11:gene-expression:protein}{प्रोटीन} संश्लेषण धीमा कर देती
है। यह संवेदनशीलता \omterm{def:g10:cells-common-unit:organelle}{सूत्रकणिकाओं} के जीवाणु पूर्वज से विरासत में मिली है।
\end{solution}

\begin{solution}{exo:g12:diversification-of-life:14}
उस \omterm{prop:g12:diversification-of-life:polyploidy}{बहुगुणित} की \omterm{def:g11:cell-cycle-mitosis:chromosome}{गुणसूत्र} संख्या अब किसी भी जनक से मेल नहीं खाती, इसलिए उनके
साथ संकरण बंध्य संतान देते हैं: अलगाव तत्काल है। जबकि किसी नए \omterm{def:g10:universal-dna:gene}{युग्मविकल्पी}
को किसी समूह को अलग दिखाने से पहले बहुत-सी पीढ़ियों तक आबादी में फैलना पड़ता
है। इसलिए कोई पादप \omterm{def:g10:biodiversity-scales:species}{जाति} एक ही पीढ़ी में उठ सकती है।
\end{solution}

\begin{solution}{exo:g12:diversification-of-life:15}
तह में, हर नया अनुक्रम --- कोई \omterm{def:g12:meiosis-genetic-shuffling:meiosis}{द्विगुणित} प्रति जो अपसरित होती है, कोई
नियामक बदलाव, किसी अंतरित प्लाज्मिड के \omterm{def:g10:universal-dna:gene}{जीन} --- कहीं न कहीं किसी \omterm{def:g11:mutations:mutation}{उत्परिवर्तन}
के रूप में शुरू हुआ। पर द्विगुणन, अंतरण और \omterm{prop:g12:diversification-of-life:polyploidy}{बहुगुणिता} पूरे-पूरे \omterm{def:g10:universal-dna:gene}{जीन} तथा
जीनोम एक ही बार में हिलाते और कई गुना करते हैं; नियामक बदलाव बिना नए
प्रोटीनों के नए रूप बनाता है; \omterm{def:g12:diversification-of-life:symbiosis}{सहजीविता} अलग-अलग \omterm{def:g10:biodiversity-scales:species}{जातियों} के जीनोम जोड़ती है;
और \omterm{prop:g12:diversification-of-life:behaviour}{संस्कृति} विविधता बिना किसी \omterm{def:g10:universal-dna:gene}{जीन} के हस्तांतरित करती है। \omterm{def:g11:mutations:mutation}{उत्परिवर्तन} अक्षर
जुटाता है; और ये प्रक्रियाएँ पूरे-पूरे पन्ने तथा किताबें दोबारा सजाती हैं।
\end{solution}

\begin{solution}{pb:g12:diversification-of-life:1}
\textbf{1.} एमर 28, रोटी वाला गेहूँ 42।

\textbf{2.} एक भी नहीं: A और B गुणसूत्र समजात नहीं हैं। युग्मकों को उन 14
में से यादृच्छिक जमावड़े मिलते हैं, पूरा सेट लगभग कभी नहीं।

\textbf{3.} 28 \omterm{def:g10:universal-dna:information}{गुणसूत्रों} के साथ हर A को एक समरूप A साथी मिलता है और हर B
को एक B: 14 जोड़े, नियमित \omterm{def:g12:meiosis-genetic-shuffling:meiosis}{अर्धसूत्री विभाजन}, 14 वाले संतुलित युग्मक।

\textbf{4.} कोई भी जोड़ी नहीं बना सकता: 7 A, 7 B और 7 D \omterm{def:g11:cell-cycle-mitosis:chromosome}{गुणसूत्र}, और हर एक
बिना किसी समजात के। बंध्य।

\textbf{5.} 21 जोड़े; और कोई परागकण 21 ढोता है।

\textbf{6.} बाक़ी दो प्रतियाँ वह \omterm{def:g11:gene-expression:protein}{प्रोटीन} जुटाती रहती हैं। अतिरिक्त प्रतियाँ
\omterm{def:g11:mutations:mutation}{उत्परिवर्तन} जमा करने और अपसरित होने को आज़ाद हैं --- यानी पूरे जीनोम के
पैमाने पर द्विगुणन।

\textbf{7.} उस इलाक़े के खेतों में जहाँ एमर उगाया जाता था, 10\,000 और
8\,000 साल पहले के बीच, जहाँ जंगली बकरी-घास 2 उस फ़सल के बीच खरपतवार की तरह
उग रही थी।

\textbf{8.} यह कि वे दो क़दम --- संकरण और द्विगुणन --- रोटी वाले गेहूँ को
उसके जनकों से बनाने के लिए काफ़ी हैं: यानी वह इतिहास दोबारा चलाया जा सकता
है।

\textbf{9.} हाँ: वह आइनकोर्न के साथ उर्वर संतान नहीं दे सकता, इसलिए उस
परिभाषा से वह एक अलग \omterm{def:g10:biodiversity-scales:species}{जाति} है, हालाँकि वह उसी से उतरा है।

\textbf{10.} बकरी-घास 2 के D \omterm{def:g11:cell-cycle-mitosis:chromosome}{गुणसूत्र} गेहूँ के D सेट के समजात हैं और उनसे
जोड़ी बनाते हैं, इसलिए कोई संकरण और उसके बाद के प्रति-संकरण वह \omterm{def:g10:universal-dna:gene}{जीन} भीतर ला
सकते हैं; जबकि किसी असंबंधित पौधे के \omterm{def:g11:cell-cycle-mitosis:chromosome}{गुणसूत्र} जोड़ी बनाते ही नहीं।

\textbf{11.} वे \omterm{def:g10:universal-dna:gene}{जीन} ख़ुद शायद ही अलग हैं; इसलिए उनकी देहों का अंतर इसी में
पड़ा होगा कि वे \omterm{def:g10:universal-dna:gene}{जीन} कैसे इस्तेमाल होते हैं।

\textbf{12.} अभिव्यक्ति के नियमन में एक बदलाव --- यानी अक्ष पर उस क्षेत्र
का विस्तार जहाँ वह \omterm{def:g10:universal-dna:gene}{जीन} सक्रिय है।

\textbf{13.} पाद वाले \omterm{def:g10:universal-dna:gene}{जीन} मौजूद हैं और काम करना शुरू करते हैं; पर वह संकेत
ग़ायब है जो पाद की बढ़त बनाए रखता, इसलिए वे कलियाँ रुक जाती हैं। फिर से
नियमन, न कि उन \omterm{def:g10:universal-dna:gene}{जीनों} का खो जाना।

\textbf{14.} यह दिखाता है कि वह \omterm{def:g10:universal-dna:gene}{जीन} कहाँ व्यक्त होता है इसे बदल देना देह के
क्षेत्र को बदलने के लिए काफ़ी है: यानी अभिव्यक्ति और रूप का सहसंबंध कारणात्मक
है।

\textbf{15.} रोटी वाले गेहूँ के बदलाव ने जीनोम में पूरे-पूरे सेट जोड़े
(उसका आकार); और अजगर के बदलाव ने यह बदला कि साझा \omterm{def:g10:universal-dna:gene}{जीन} कहाँ इस्तेमाल होते हैं
(उनका नियमन)।

\textbf{16.} अपना वृत्ताकार \omterm{def:g10:universal-dna:information}{DNA} (अपेक्षित: कोई नहीं, सारे \omterm{def:g10:universal-dna:gene}{जीन} केंद्रकीय);
जीवाणु \omterm{def:g10:cells-common-unit:organelle}{राइबोसोम} (अपेक्षित: सुकेंद्रकी \omterm{def:g10:cells-common-unit:organelle}{राइबोसोम}); विभाजन से गुणन (अपेक्षित:
पुर्ज़ों से जोड़ाई); दोहरी झिल्ली (अपेक्षित: बाक़ी \omterm{def:g10:cells-common-unit:organelle}{कोशिकांगों} जैसी एक ही
झिल्ली)। और उनके \omterm{def:g10:universal-dna:gene}{जीन} जीवाणु \omterm{def:g10:universal-dna:gene}{जीनों} जैसे हैं (अपेक्षित: \omterm{def:g10:cells-common-unit:organelle}{केंद्रक} प्रकार के
\omterm{def:g10:universal-dna:gene}{जीन})।

\textbf{17.} \omterm{def:g10:cells-common-unit:organelle}{केंद्रक} में: समय के साथ उस जीवाणु के अधिकांश \omterm{def:g10:universal-dna:gene}{जीन} केंद्रकीय
जीनोम में अंतरित हो गए, और उस \omterm{def:g10:cells-common-unit:organelle}{कोशिकांग} का अपना जीनोम सिकुड़कर एक अवशेष रह
गया।

\textbf{18.} \omterm{def:g10:cells-common-unit:organelle}{सूत्रकणिका} \omterm{def:g10:universal-dna:gene}{जीन} केवल माँ से संतान तक जाते हैं; और किसी व्यक्ति
का \omterm{def:g10:cells-common-unit:organelle}{सूत्रकणिका} \omterm{def:g10:universal-dna:information}{DNA} अटूट मातृ वंश का पीछा करता है।

\textbf{19.} \omterm{def:g10:cells-common-unit:organelle}{सूत्रकणिका} ने अपने अधिकांश \omterm{def:g10:universal-dna:gene}{जीन} अपने पोषक को सौंप दिए हैं और वह
\omterm{def:g10:cells-common-unit:cell}{कोशिका} के बाहर रह ही नहीं सकती; जबकि लाइकेन के साझीदार अपने जीनोम रखते हैं
और मुश्किल से ही सही, अलग किए जा सकते हैं।

\textbf{20.} आइनकोर्न (14, A), बकरी-घास 1 (14, B), बकरी-घास 2 (14, D):
42 \omterm{def:g11:cell-cycle-mitosis:chromosome}{गुणसूत्र}; ``नियमन''; और वह \omterm{def:g10:cells-common-unit:organelle}{सूत्रकणिका}।
\end{solution}
